\documentclass[10pt]{article}
\usepackage{amsmath}
\usepackage{amssymb}
\usepackage{hyperref}
\usepackage[svgnames]{xcolor} 
\usepackage{listings}
\usepackage[skins,listings,breakable]{tcolorbox}

\definecolor{cppbg}{HTML}{F8F9FA}     
\definecolor{boxtitlebg}{HTML}{2B2D31}
\definecolor{cppblue}{HTML}{00599C}   
\definecolor{cppgreen}{HTML}{138A07} 
\definecolor{cppred}{HTML}{B11515}  
\definecolor{cppnumber}{HTML}{888888} 

\lstdefinestyle{cppstyle}{
    language=C++,
    commentstyle=\color{cppgreen}\itshape,
    keywordstyle=\color{cppblue}\bfseries,
    numberstyle=\tiny\color{cppnumber},
stringstyle=\color{cppred},
basicstyle=\ttfamily\footnotesize,
breaklines=true,
    showstringspaces=false,
    tabsize=4,                    
    numbers=none                     
}

\newtcblisting{cppbox}[1]{
    enhanced,
    skin=bicolor,
    colback=cppbg,
    colbacktitle=boxtitlebg,
    coltitle=white,
    colframe=boxtitlebg,
    boxrule=0.5pt,
    arc=4pt,
    drop shadow=black!10!white,
    fonttitle=\bfseries\sffamily,
    title=#1,
    listing only,
    listing options={style=cppstyle, numbers=left, numbersep=10pt},
    left=15pt,
    top=4pt,
    bottom=4pt,
    breakable
}



\title{E.A.6.5 - Schur's lemma}
\author{}
\date{}

\begin{document}
\maketitle
\vspace{-6em}

\section{Modellazione}

Definiamo le variabili atomiche $X_{i,j}$ t.c. $X_{i,j} = \top \iff$ la biglia $i \in \{1, \dots, n\}$ si trova nell'urna $j \in \{1, 2, 3\}$

La formula proposizionale è la congiunzione dei seguenti vincoli:

\begin{enumerate}
    \item ogni biglia $i$ deve essere assegnata ad almeno un'urna
        $$ \varphi_{ex} = \bigwedge_{i=1}^{n} (X_{i,1} \lor X_{i,2} \lor X_{i,3}) $$

    \item ogni biglia non può trovarsi in più di un'urna
    $$ \varphi_{un} = \bigwedge_{i=1}^{n} \left( (\neg X_{i,1} \lor \neg X_{i,2}) \land (\neg X_{i,1} \lor \neg X_{i,3}) \land (\neg X_{i,2} \lor \neg X_{i,3}) \right) $$

    \item per ogni urna $j$, non devono esistere tre biglie $x, y, z$ tali che $x + y = z$    
        $$ \varphi_{s} = \bigwedge_{j=1}^{3} \bigwedge_{\substack{1 \le x, y, z \le n \\ x+y=z}} (\neg X_{x,j} \lor \neg X_{y,j} \lor \neg X_{z,j}) $$
        \subitem (consideriamo anche il caso in cui le biglie non siano distinte (es. 1+1=2))
\end{enumerate}


Ovvero
$$ \Phi_n = \varphi_{ex} \land \varphi_{un} \land \varphi_{s} $$

\section{Istanziazione}

\begin{enumerate}
    \item vincoli di esistenza e unicità
\begin{itemize}
    \item \textbf{biglia 1}: $(X_{1,1} \lor X_{1,2} \lor X_{1,3}) \land (\neg X_{1,1} \lor \neg X_{1,2}) \land (\neg X_{1,1} \lor \neg X_{1,3}) \land (\neg X_{1,2} \lor \neg X_{1,3})$
    \item \textbf{biglia 2}: $(X_{2,1} \lor X_{2,2} \lor X_{2,3}) \land (\neg X_{2,1} \lor \neg X_{2,2}) \land (\neg X_{2,1} \lor \neg X_{2,3}) \land (\neg X_{2,2} \lor \neg X_{2,3})$
    \item \textbf{biglia 3}: $(X_{3,1} \lor X_{3,2} \lor X_{3,3}) \land (\neg X_{3,1} \lor \neg X_{3,2}) \land (\neg X_{3,1} \lor \neg X_{3,3}) \land (\neg X_{3,2} \lor \neg X_{3,3})$
    \item \textbf{biglia 4}: $(X_{4,1} \lor X_{4,2} \lor X_{4,3}) \land (\neg X_{4,1} \lor \neg X_{4,2}) \land (\neg X_{4,1} \lor \neg X_{4,3}) \land (\neg X_{4,2} \lor \neg X_{4,3})$
    \item \textbf{biglia 5}: $(X_{5,1} \lor X_{5,2} \lor X_{5,3}) \land (\neg X_{5,1} \lor \neg X_{5,2}) \land (\neg X_{5,1} \lor \neg X_{5,3}) \land (\neg X_{5,2} \lor \neg X_{5,3})$
\end{itemize}

    \item vincoli di Schur
        \subitem le terne $(x, y, z)$ tali che $x+y=z$ sono: 
        \subitem $(1,1,2), (1,2,3), (1,3,4), (1,4,5), (2,2,4), (2,3,5)$

    \begin{itemize}
    \item \textbf{terna (1,1,2)}: $(\neg X_{1,1} \lor \neg X_{2,1})$
    \item \textbf{terna (1,2,3)}: $(\neg X_{1,1} \lor \neg X_{2,1} \lor \neg X_{3,1})$
    \item \textbf{terna (1,3,4)}: $(\neg X_{1,1} \lor \neg X_{3,1} \lor \neg X_{4,1})$
    \item \textbf{terna (1,4,5)}: $(\neg X_{1,1} \lor \neg X_{4,1} \lor \neg X_{5,1})$
    \item \textbf{terna (2,2,4)}: $(\neg X_{2,1} \lor \neg X_{4,1})$
    \item \textbf{terna (2,3,5)}: $(\neg X_{2,1} \lor \neg X_{3,1} \lor \neg X_{5,1})$
\end{itemize}

\begin{itemize}
    \item \textbf{terna (1,1,2)}: $(\neg X_{1,2} \lor \neg X_{2,2})$
    \item \textbf{terna (1,2,3)}: $(\neg X_{1,2} \lor \neg X_{2,2} \lor \neg X_{3,2})$
    \item \textbf{terna (1,3,4)}: $(\neg X_{1,2} \lor \neg X_{3,2} \lor \neg X_{4,2})$
    \item \textbf{terna (1,4,5)}: $(\neg X_{1,2} \lor \neg X_{4,2} \lor \neg X_{5,2})$    
    \item \textbf{terna (2,2,4)}: $(\neg X_{2,2} \lor \neg X_{4,2})$
    \item \textbf{terna (2,3,5)}: $(\neg X_{2,2} \lor \neg X_{3,2} \lor \neg X_{5,2})$
\end{itemize}

\begin{itemize}
    \item \textbf{terna (1,1,2)}: $(\neg X_{1,3} \lor \neg X_{2,3})$
    \item \textbf{terna (1,2,3)}: $(\neg X_{1,3} \lor \neg X_{2,3} \lor \neg X_{3,3})$
    \item \textbf{terna (1,3,4)}: $(\neg X_{1,3} \lor \neg X_{3,3} \lor \neg X_{4,3})$
    \item \textbf{terna (1,4,5)}: $(\neg X_{1,3} \lor \neg X_{4,3} \lor \neg X_{5,3})$
    \item \textbf{terna (2,2,4)}: $(\neg X_{2,3} \lor \neg X_{4,3})$    
    \item \textbf{terna (2,3,5)}: $(\neg X_{2,3} \lor \neg X_{3,3} \lor \neg X_{5,3})$
\end{itemize}

\end{enumerate}

\section{Codec}
(encoder in cpp)

\begin{cppbox}{}
#pragma once

#include "Encoder.hpp"
#include "Problem.hpp"
#include <string>
#include <vector>

class SchursLemma : public Problem {
private:
    int N;

    inline std::string vName(int ball, int urn) const {
        return "X_" + std::to_string(ball) + "_" + std::to_string(urn);
    }

public:
    inline SchursLemma(int n) : N(n) {}

    void generateConstraints(Encoder& encoder, EncodingContext& ctx, int tId, int tCount) override {
        
        // esistenza e unicita'
        // (ogni thread gestisce le biglie i tali che (i % tCount == tId))
        for (int i = 1 + tId; i <= N; i += tCount) {
            std::vector<int> lits;
            for (int j = 1; j <= 3; ++j) {
                lits.push_back(encoder.getVar(vName(i, j)));
            }

            // ogni biglia in almeno un'urna
            encoder.writeDirectClause(lits, ctx);

            // ogni biglia in al massimo un'urna
            encoder.writeDirectClause({-lits[0], -lits[1]}, ctx);
            encoder.writeDirectClause({-lits[0], -lits[2]}, ctx);
            encoder.writeDirectClause({-lits[1], -lits[2]}, ctx);
        }

        // vincolo di schur
        for (int x = 1 + tId; x <= N; x += tCount) {
            // assumiamo x <= y per evitare di generare due volte la
            // stessa combinazione (es. 1+2=3 e 2+1=3 generano le stesse clausole)
            for (int y = x; x + y <= N; ++y) {
                int z = x + y;
                
                for (int j = 1; j <= 3; ++j) {
                    int vx = encoder.getVar(vName(x, j));
                    int vy = encoder.getVar(vName(y, j));
                    int vz = encoder.getVar(vName(z, j));

                    encoder.writeDirectClause({-vx, -vy, -vz}, ctx);
                }
            }
        }
    }
};
\end{cppbox}



\end{document}
